\documentclass[11pt]{article}

\usepackage{amsmath,amssymb,amsthm}
\usepackage{graphicx}
\usepackage{hyperref}
\usepackage{booktabs}
\usepackage{enumitem}
\usepackage[margin=1in]{geometry}

\newtheorem{theorem}{Theorem}
\newtheorem{remark}{Remark}

\title{When Meta-Learned Optimization Fails to Beat Round-Robin:\\
A Validation Study on Hyperbolic Polytope Volume}

\author{Igor Rivin}

\date{\today}

\begin{document}

\maketitle

\begin{center}
\fbox{\parbox{0.92\linewidth}{\textbf{Superseded validation result (2026-08-15).}
The evaluation-accounting audit found that the baseline runner charged padded
history lengths rather than actual objective, surrogate, and gradient work.
The ``matched budget'' claims and numerical comparisons below are therefore
not valid current evidence and must not be cited or published. The source is
retained for provenance; rerun the study with the repaired validator and
retrained controllers before replacing this notice.}}
\end{center}

\begin{abstract}
We investigate whether a reinforcement-learning meta-controller can
outperform fixed single-optimizer and uniform-mixing baselines on a
family of optimization tasks parameterized by dimension $d$ and
vertex count $N$ (maximizing the hyperbolic volume of an ideal
polytope in $\mathbb{H}^d$).  A Gymnasium environment with a
hierarchical search tree dispatches calls to a small library of inner
optimizers (CMA-ES on the Euclidean surrogate, L-BFGS on closed-form
hyperbolic gradients, plus generic sample / split / perturb / prune
operations); a PPO policy is trained per dimension for
$200{,}000$ environment steps.  Evaluated at matched budget on
$200$ tasks across $d \in \{3, 4, 5, 6\}$, the trained controller
wins cleanly only at $d = 5$ ($+15.9\%$ geometric-mean ratio over
the best single baseline).  At $d = 3$ and $d = 4$ it loses to a
round-robin baseline that splits the budget equally across the same
inner optimizers; at $d = 6$ it is dominated by a single specialized
solver (L-BFGS on the Murakami $\mathbb{H}^6$ volume), with
catastrophic failure at the simplex-boundary case $N = 7$.  We
report the negative result honestly: meta-optimization adds value
where no single classical heuristic dominates, and round-robin
suffices everywhere else.  All code, saved policies, and validation
harness are released so the reward and observation design can be
iterated on by other researchers.
\end{abstract}

% ============================================================================
\section{Setup}
% ============================================================================

The volume-maximization problem family parameterized by $(d, N)$
is a natural target for meta-learned optimization: each task asks
the same question (find a high-volume ideal polytope) with a shared
structure (points on a sphere, a few closed-form objective and
gradient routines, a handful of useful inner optimizers), and
optimizer choice is clearly non-uniform across tasks---CMA-ES on
the Euclidean surrogate dominates at small $N$, L-BFGS on the
hyperbolic objective dominates once the configuration is close
to optimal, and local sampling plus region splitting help
escape basins at larger $N$.  We implemented a hierarchical
reinforcement-learning harness to test whether a learned
meta-controller can outperform fixed single-optimizer and
uniform-mixing baselines at matched evaluation budget.

\paragraph{Harness.}
The harness is a Gymnasium environment wrapping a hierarchical
search tree.  Each state holds the frontier of active regions of
the configuration space $(S^{d-1})^N$; the agent selects a region
and an operation (sample, split, perturb, prune, or invoke one of
the inner optimizers).  The reward is the improvement in the
best-so-far volume minus a small cost per function evaluation.
A PPO policy (stable-baselines3, MLP backbone) is trained per
dimension for $200{,}000$ environment steps, which corresponds
to roughly $40$--$400$ hours of simulated optimization depending
on the inner-optimizer cost.  The same harness supports problems
orthogonal to polytope volume (we use it for cubic-graph girth
maximization), which is the sense in which it is problem-agnostic.

\paragraph{Inner optimizers.}
For the polytope task family the inner-optimizer slate is:
\begin{itemize}
\item \textbf{CMA-ES on Euclidean surrogate:} fast, low-cost per
      evaluation; targets the convex-hull volume of the inscribed
      polytope.  Effective at small $N$ where the optimum is
      algebraically trivial.
\item \textbf{L-BFGS on the hyperbolic objective:} closed-form
      gradients via Bloch--Wigner ($d = 3$), the Haagerup--Munkholm
      formula ($d = 4$), the Schl\"afli surrogate ($d = 5$), or
      the Murakami $\mathbb{H}^6$ volume ($d = 6$).  Slow per call
      but converges to high-precision local maxima.
\item \textbf{Sample, split, perturb, prune:} generic
      tree-management actions.
\end{itemize}

\paragraph{Fix to a latent scoring bug.}
The pre-refactor environment tracked a mix of Euclidean and
hyperbolic volumes in the same node-best field at $d \in
\{3, 4, 6\}$, because the CMA path scored Euclidean while the
L-BFGS path returned hyperbolic.  In the refactored harness used
for the present study, each task has one consistent dimension-specific
scoring function (closed-form hyperbolic for $d \in \{3, 4, 6\}$;
Euclidean at $d = 5$, where it agrees with the hyperbolic optimum
to full precision).  All baseline numbers below are reported under
the corrected scoring.

% ============================================================================
\section{Validation Grid and Baselines}
% ============================================================================

We test on a $4 \times 5 \times 10$ grid: $d \in \{3, 4, 5, 6\}$,
five $N$ values per $d$ chosen to span small $N$ where the global
optimum is algebraically trivial (the regular simplex at $N = d+1$,
the 4-simplex bipyramid, etc.) and larger $N$ where combinatorial
variety is rich, and $10$ independent seeds per $(d, N)$.  Three
methods are compared at matched budget (the evaluation count
consumed by one RL episode):
\begin{enumerate}[label=(\roman*)]
\item \textbf{RL:} the trained PPO controller;
\item \textbf{Round-robin:} the same total budget split equally
      across the inner optimizers available to the RL agent,
      ignoring context;
\item \textbf{Solo:} each single inner optimizer given the whole
      budget; we report the best single solo per task.
\end{enumerate}
The pre-registered decision rule in our protocol is: RL passes the
\emph{learning-allocation} test if and only if it beats round-robin,
because round-robin already exposes the same diversity of inner
optimizers without any learning.

% ============================================================================
\section{Results}
% ============================================================================

The outcome splits sharply by dimension
(Table~\ref{tab:rl_per_dim}).

\begin{table}[ht]
\centering
\caption{Geometric-mean ratio to the best single solo baseline,
per dimension, over 50 tasks ($5$ $N$'s $\times$ $10$ seeds).
Values $> 1$ mean the method beats the strongest fixed optimizer.}
\label{tab:rl_per_dim}
\smallskip
\begin{tabular}{r c c l}
\toprule
$d$ & RL & Round-robin & Dominant single-solo baseline \\
\midrule
3 & 1.057 & \textbf{1.081} & L-BFGS on Bloch--Wigner \\
4 & 1.132 & \textbf{1.238} & CMA-ES on Euclidean surrogate \\
5 & \textbf{1.159} & 1.049 & CMA-ES on Euclidean surrogate \\
6 & 0.580 & 0.529 & \textbf{L-BFGS on Murakami $\mathbb{H}^6$} (ratio $0.980$) \\
\bottomrule
\end{tabular}
\end{table}

The RL controller wins cleanly only at $d = 5$.  At $d = 3$ and
$d = 4$ the uninformed round-robin baseline---splitting the
budget equally across inner optimizers---beats the trained PPO
policy by $2$--$9\%$.  At $d = 6$ a single specialized
inner optimizer (L-BFGS on the Murakami hyperbolic volume) is
dominant: both RL and round-robin are far below it.

\paragraph{The simplex-boundary failure at $d = 6, N = 7$.}
The $d = 6$ deficit is driven primarily by the smallest $N$ in the
grid.  At $N = 7$ (the regular ideal $6$-simplex), the solo L-BFGS
baseline finds the global optimum $V \approx 0.0067$ in a few
gradient steps, while the RL median is $0.0005$---a factor of $13$
below.  For $N \in \{10, 16, 20, 24\}$ at $d = 6$ the RL controller
ties the solo-L-BFGS baseline to three significant digits,
indicating that the learned policy \emph{can} identify the right
specialized solver in the large-$N$ regime but does not recognize
the simplex-boundary special case at small $N$.  The pattern
reproduces, less dramatically, at $d \in \{3, 4, 5\}$:

\begin{center}
\begin{tabular}{r r r r}
\toprule
$d$ & $N_{\min}$ (simplex) & RL/solo at $N_{\min}$ & RL/solo for largest $N$ \\
\midrule
3 & 5  & 1.000 & 1.081 \\
4 & 5  & 0.977 & 1.192 \\
5 & 6  & 0.905 & 1.287 \\
6 & 7  & \textbf{0.079} & 1.001 \\
\bottomrule
\end{tabular}
\end{center}

The penalty for RL over-exploration grows sharply with $d$ at the
simplex boundary, suggesting the policy has not learned to
``defer'' to the dominant inner optimizer when the configuration
is close to a known closed-form optimum.

% ============================================================================
\section{Narrow Claim and Failure Modes}
% ============================================================================

The usable takeaway is more modest than ``RL beats classical
optimization on hyperbolic polytope volume'':

\begin{enumerate}
\item \textbf{Where RL adds value:} at $d = 5$, the
      Euclidean--hyperbolic volume equivalence is exact (no inner
      optimizer dominates uniformly across $N$), and the learned
      controller captures enough of the right heuristic
      ($+10.5\%$ over round-robin, $+15.9\%$ over best solo
      baseline) to outperform uninformed mixing.

\item \textbf{Where RL ties or loses:} at $d \in \{3, 4\}$ a
      strong solo baseline (L-BFGS on Bloch--Wigner; CMA-ES on
      Euclidean) suffices, and round-robin already exposes the
      diversity that RL would need to learn.  The trained policy
      contributes nothing beyond its noise floor.

\item \textbf{Where RL fails badly:} at $d = 6$ the closed-form
      Murakami gradient is dominant; the simplex-boundary case
      $N = 7$ exposes a specific failure mode in the learned
      policy (over-exploration when convergence is needed).
\end{enumerate}

This is consistent with the broader meta-learning literature:
meta-optimization pays off where classical heuristics are weak;
where a strong specialized solver exists, the learned policy
has to discover that solver's dominance and defer to it, which
the current policy does only partially.

\paragraph{Why round-robin is so strong.}
The strength of round-robin in our setting is itself an
informative finding.  The inner-optimizer slate is small (4 to 6
options), each option is itself reasonable, and the per-call cost
varies by $\sim 10\times$ across optimizers.  Allocating the
budget uniformly across them implicitly hedges across regimes:
when CMA-ES dominates (small $N$, easy landscape) it gets a
quarter of the budget and that suffices; when L-BFGS dominates
(close to optimum) it likewise gets enough.  A learned policy
that is no better than uniform allocation is failing to add the
context-dependent allocation that justifies its training cost.

% ============================================================================
\section{Reproducibility}
% ============================================================================

The full training scripts, saved policies, and validation harness
are released at the repository accompanying this note.  The
training command is
\begin{verbatim}
python scripts/train_polytope_controller.py --d <D> --steps 200000
\end{verbatim}
producing \texttt{ppo\_d<D>\_v2\_final.zip} plus a configuration
sidecar.  The validation grid is reproduced by
\begin{verbatim}
python scripts/validate_rl_controller.py \
    --model-template 'results/rl_models/ppo_d{d}_v2_final.zip' \
    --dims 3,4,5,6 --seeds 10
\end{verbatim}
which writes a JSON of every $(\text{method}, d, N, \text{seed})$
result, plus the aggregate ratios reported above.  Both scripts
share the same harness, so a researcher who wants to test a
different reward function, observation encoding, or inner-optimizer
slate can change them in one place and rerun the full grid in
roughly a day of CPU time on a modest workstation.

% ============================================================================
\section{Future Work}
% ============================================================================

Three directions seem worth pursuing:

\begin{enumerate}
\item \textbf{Reward shaping.} The current reward is the increment
      in best-so-far volume.  A potential-based shaping that
      includes a constant per-step penalty for over-exploration at
      the simplex boundary, or a curriculum that starts at large
      $N$ where RL has its advantage, might rescue the small-$N$
      cases.

\item \textbf{Observation features.} The current observation is
      a flat vector of per-region statistics (depth, mean, variance,
      best, etc.) plus task metadata $(d, N, \log Nd)$.  Adding a
      symmetry-invariant fingerprint of the current best
      configuration---e.g., the multiset of pairwise chordal
      distances---would let the policy recognize ``this is close
      to a regular simplex'' and defer to the corresponding solo
      optimizer.

\item \textbf{Comparison against off-the-shelf hyperparameter
      tuners.} Optuna and similar tools target a related but not
      identical task (single-study tuning, not learned meta-control
      over a task family).  Benchmarking against them on the same
      $(d, N)$ grid would establish whether our negative result is
      specific to our PPO setup or generalizes across meta-learning
      approaches.
\end{enumerate}

\end{document}
